\subsection{Integral theta lifts and exact canonical section norms}

\begin{theorem}[DP62: individual allowed translates]\label{DP62}
Use the good model of the Jacobian, its symmetric rigidified trace-one line $L$, and its principal polarization over
$S=\operatorname{Spec}\mathcal O_E[1/14]$, where $E=\mathbf Q(\zeta_7)$.
Fix a nonzero algebraic section $s\in H^0(A,L)$ primitive at the selected prime over $11$.
Let $M$ be odd, put $N=M^2$, and let
\[
 \Phi_M:[M]^*L\xrightarrow{\sim}L^N
\]
be the canonical rigidified multiplication isomorphism.
For a finite extension $F/E$ and $t\in A[N](F)$, the theta lift $U_t$ with
\[
 U_t^N=1,\qquad \iota(U_t)=U_t^{-1}
\]
exists uniquely over the field of definition of $t$ and extends integrally over $\mathcal O_F[1/14]$.
This includes places dividing $M$.
For any local lift $U$, it is given by
\begin{equation}\label{eq:rootfreethetalift}
 c=\iota(U)U,\quad d=U^N,\qquad
 U_t=c^{(N-1)/2}d^{-1}U.
\end{equation}
The formula glues without any root extraction.
Consequently every individual translate
\[
 s_t=U_t\bigl(\Phi_M([M]^*s)\bigr)
\]
is integral on the good model over $\mathcal O_F[1/14]$.

At every complex embedding, equip $L$ with the canonical smooth Hermitian metric whose curvature is translation-invariant and whose rigidification at zero is an isometry.
With the tensor metric on $L^N$, one has the exact equality
\begin{equation}\label{eq:individualthetanorm}
 \|s_t\|_{\sigma,\sup}
 =\|s\|_{\sigma|_E,\sup}.
\end{equation}
The same holds for $L^r$-norms, $1\le r<\infty$, with probability Haar measure.

For the permitted subgroup $H=A[\alpha^2]$ of DP59, any $r$-tuple of points, its canonical lifts and its translated sections are defined over a field of degree at most
\[
 [E(t_1,\ldots,t_r):E]\le M^{2gr}.
\]
Thus the ordinary and two-copy harmonic sharp witnesses in dimension three need degrees at most $M^6$ and $M^{12}$.
Their canonical section norms are uniformly bounded, and they have no new good-prime denominators, at the levels $M=29^e=O(n)$ and $M^2=O(n^2)$ of DP59.
No discriminant or complete source derivative/gauge-cost bound is asserted.
\end{theorem}

\begin{proof}
\emph{The finite-flat theta torsor.}
For the principally polarized line,
\[
 \phi_{L^N}=N\phi_L,\qquad K(L^N)=A[N].
\]
This kernel is finite locally free even at primes dividing $N$.
Indeed each fibre of multiplication by $N$ is finite locally free of degree $N^{2g}$ by the known field theorem, Stacks Project Lemma~39.9.8, Tag~0BFG.
The relative source is smooth, hence flat and locally of finite presentation over $S$; the target is locally of finite type.
The fibrewise criterion, Lemma~37.16.4, Tag~039E, gives relative flatness.
The map is proper with finite fibres and hence finite by Lemma~37.44.1, Tag~02LS.
Over this noetherian base, finite flatness gives local freeness, of the stated fibre rank.
The kernel is the base change along zero.\footnote{The exact known statements are \url{https://stacks.math.columbia.edu/tag/0BFG}, \url{https://stacks.math.columbia.edu/tag/039E}, and \url{https://stacks.math.columbia.edu/tag/02LS}.}

The relative theta group is the central extension
\[
 1\longrightarrow\mathbf G_m\longrightarrow
 \mathcal G(L^N)\longrightarrow A[N]\longrightarrow1.
\]
Its fibre is the torsor of translation isomorphisms of $L^N$.
The normalized symmetry induces an automorphism $\iota$ which fixes the central scalars and covers inversion on $A[N]$.
The relative finite-flat setting and this symmetry convention are the known theta-group statements in Polishchuk, \emph{Analogue of Weil representation for abelian schemes}, \S1, author preprint pp.~4--5, before Lemma~1.1.
Proposition~1.2 and Corollary~1.3 also establish the odd symmetric torsor's triviality; the following formula makes its normalization explicit.\footnote{\url{https://arxiv.org/pdf/alg-geom/9712021}. The stated page numbers refer to this directly read author version.}

\emph{Normalization over arbitrary test bases.}
For a point $t\in A[N](T)$, choose a lift $U$ locally in the fppf topology.
The elements $d=U^N$ and $c=\iota(U)U$ lie over zero, so are invertible central scalars.
Since $\iota(U)=cU^{-1}$, these two lifts commute, and
\[
 c^N=\iota(U)^NU^N=d^2.
\]
For $a=c^{(N-1)/2}d^{-1}$ this gives
\[
 a^2c=c^Nd^{-2}=1,\qquad
 a^Nd=c^{N(N-1)/2}d^{1-N}=1.
\]
Hence $aU$ has both required identities.
Replacing the local lift by $uU$ changes $c$ to $u^2c$, $d$ to $u^Nd$, and $a$ to $u^{-1}a$.
Thus $aU$ is unchanged and descends on every overlap.
Only multiplication and inversion of units occur; neither $N$ nor any root is inverted or adjoined.
The construction commutes with arbitrary base change.
Two normalized lifts differ by a unit with $u^2=u^N=1$; since $N$ is odd, an integer B\'ezout identity gives $u=1$, including over nonreduced bases.

This is a section of the underlying torsor, not a group splitting of the whole central extension.
On an isotropic subgroup the lifts commute; their products again satisfy the two identities, so uniqueness makes that restriction a homomorphism.
The normalized symmetry and the power identity are exactly the clauses of the recovered \texttt{ThetaPt.SymLift} definition, lines~44--51; the action on sections is \texttt{IsActOf}, lines~83--84.

\emph{Integrality of the individual algebraic sections.}
A torsion point over $F$ extends uniquely over every good DVR by properness of the abelian scheme.
Stacks Project Lemma~29.43.5, Tag~0BX7, extends the rational map to a neighborhood of each such point: the source is noetherian and its local ring is a valuation ring.
The valuative criterion is Lemma~29.43.1, Tag~0BX5.
The neighborhood extensions agree on overlaps by separatedness and the dense common generic fibre, so they glue over $\operatorname{Spec}\mathcal O_F[1/14]$.\footnote{\url{https://stacks.math.columbia.edu/tag/0BX7}; \url{https://stacks.math.columbia.edu/tag/0BX5}.}
Separatedness preserves $Nt=0$.
Applying \eqref{eq:rootfreethetalift} to the resulting integral theta torsor gives an integral lift on this same base.
In particular the field of the point already contains the canonical lift.

DP59's Poincar\'e construction makes the symmetry and $\Phi_M$ integral on the good model.
Pulling back the fixed integral section by $[M]$, applying $\Phi_M$, and then applying the integral theta action therefore gives $s_t$ integrally.
This includes the places above $29$ when $M=29^e$; their torsion kernels need not be \'etale.
For these DP59 levels the torsion at $11$ is finite \'etale, so the already unramified point field requires no further extension for the lift.
The integer valuation group used in DP59 is preserved.
No assertion that the fixed section is primitive at every other good prime is needed.
The existence of a fixed global $s$ primitive at $11$ follows by clearing the local generator's denominators outside that prime, as in DP60.

\emph{Canonical metrics.}
The known metric theorem is Moret-Bailly, \emph{Pinceaux de vari\'et\'es ab\'eliennes}, Ast\'erisque~129 (1985), Proposition~II.2.1, printed p.~48: a cubical line on a complex abelian variety has its unique smooth Hermitian metric compatible with the rigidification and with translation-invariant curvature.
The equivalent characterization is also stated in de Jong--Shokrieh, \emph{Faltings height and N\'eron--Tate height of a theta divisor}, Compositio Mathematica~158 (2022), \S12.2.\footnote{Original primary page: \url{https://www.numdam.org/item/AST_1985__129__1_0.pdf}, printed p.~48, PDF page~51. The original page and the later published statement were both directly read.}
We use this known theorem, rather than re-proving metric existence.

Let $\omega$ be the invariant curvature of $L$ at a fixed complex embedding.
Because the differential of $[M]$ is multiplication by $M$,
\[
 [M]^*\omega=M^2\omega=N\omega.
\]
The pullback and tensor metrics compared by $\Phi_M$ therefore have the same curvature.
Their logarithmic ratio is a smooth real pluriharmonic function on the compact connected torus and hence is constant.
The rigidified isomorphism is an isometry at zero, so that constant is zero.
Thus $\Phi_M$ is isometric and
\[
 \|\Phi_M([M]^*s)(x)\|=\|s([M]x)\|.
\]
Surjectivity of $[M]$ gives equality of the sup norms.

For a finite-order theta lift $U_t$, the comparison of the metric with its translate has zero curvature because $T_t^*(N\omega)=N\omega$.
Its positive norm ratio is therefore a constant $r$.
Composing the lift $N$ times gives ratio $r^N$, while $U_t^N=1$ gives ratio one.
Hence $r=1$: the theta lift is isometric.
Translation is bijective, proving \eqref{eq:individualthetanorm}.
With probability Haar measure, translation preserves measure and the pushforward by $[M]$ is again probability Haar measure; this proves the finite $L^r$ equalities.
Every embedding of $F$ restricts to one of the finitely many embeddings of the fixed field $E$, so all the constants are uniform in the level and the point.

\emph{The actual tuple field.}
The subgroup $H=A[\alpha^2]$ is defined over $E$, is \'etale in characteristic zero, and has degree $M^{2g}$ by DP59.
Consequently $H^r$ has degree $M^{2gr}$, bounding the residue-field degree of any tuple by that number.
The normalized lift is unique and functorial, so its coefficients and the translated section descend to the same field.
This gives the stated $M^6$ and $M^{12}$ bounds in dimension three.
For comparison, the full torsion splitting field has Galois group injecting into $\operatorname{GL}_{2g}(\mathbf Z/N\mathbf Z)$ and therefore degree at most $N^{4g^2}=M^{8g^2}$, though no full splitting field is needed for an individual tuple.
No discriminant estimate follows merely from either degree bound.
\end{proof}

The norm in DP62 is the canonical norm of the section itself.
Dividing by a source analytic frame, multiplying by its exponential gauge, and applying a specified differential polynomial can introduce additional level and degree factors.
Those operations, bad-prime norm comparisons, discriminants, prescribed row weights, the mixed-type moduli interface and membership in the iterated theta image require their own estimates.
The theorem removes an uncontrolled canonical section norm and new good-prime denominators for the individual allowed witnesses; it does not by itself compare the net cost of whole theta-lift chains.
